\documentclass[12pt,a4paper]{article}
\usepackage[utf8]{inputenc}
\usepackage{geometry}
\geometry{margin=1in}
\usepackage{hyperref}
\usepackage{booktabs}
\usepackage{listings}
\usepackage{xcolor}

\lstset{
    basicstyle=\ttfamily\footnotesize,
    breaklines=true,
    frame=single,
    backgroundcolor=\color{gray!10},
    keepspaces=true
}

\title{\textbf{Cryptography and Network Security Technical Report}\\ \large Integrated Situation Assessment}
\author{Registration Number: 4202670014}
\date{\today}

\begin{document}

\maketitle

\section{Risk Assessment}
A security review of the polytechnic institute identified critical vulnerabilities across core assets. Table \ref{tab:risks} summarizes these assets, vulnerabilities, consequences, risk rankings, and recommended controls.

\begin{table}[h!]
\centering
\small
\begin{tabular}{@{}lllll@{}}
\toprule
\textbf{Asset} & \textbf{Vulnerability} & \textbf{Consequence} & \textbf{Rank} & \textbf{Control} \\ \midrule
Records Server & Guest access allowed & Data breach / leak & Critical & Firewall / ACL Segmentation \\
Data in Transit & Unencrypted transfers & Eavesdropping / MitM & High & AES Encryption / SFTP \\
Staff Accounts & Weak passwords / OS & System compromise & Medium & Password Policy \& Patching \\ \bottomrule
\end{tabular}
\caption{Risk Assessment Matrix}
\label{tab:risks}
\end{table}

\section{Encryption and Integrity Application}
To secure sensitive student records, a Python security toolkit (\texttt{crypto\_toolkit.py}) was developed using the \texttt{cryptography} library.

\subsection{Implementation Details}
\begin{itemize}
    \item \textbf{Encryption \& Decryption:} Implemented via symmetric encryption using Fernet (AES-128 in CBC mode with HMAC authentication). The key is generated and stored outside the repository directory (\texttt{../secret.key}).
    \item \textbf{Integrity Checking:} Calculates and verifies SHA-256 hashes before encryption and after decryption to confirm the restored file is identical to the original.
    \item \textbf{Error Handling:} Uses standard Python \texttt{try-except} blocks to catch \texttt{FileNotFoundError} and \texttt{InvalidToken} exceptions gracefully without terminating execution.
\end{itemize}

\subsection{Execution and Test Evidence}
Below is the terminal output from running \texttt{crypto\_toolkit.py}. Note that the output lines starting with \texttt{[-] Error:...} under Section 4 are **expected test outputs** demonstrating that the error-handling functions successfully capture missing file errors and corrupted decryption attempts without crashing the program:

\begin{lstlisting}[caption={Python Security Toolkit Execution Output}]
[+] New key generated and saved to '../secret.key'.

--- 1. ENCRYPTION ---
[+] File successfully encrypted and saved to 'sample_student_records.csv.enc'.

--- 2. DECRYPTION ---
[+] File successfully decrypted and saved to 'sample_student_records_decrypted.csv'.

--- 3. INTEGRITY CHECK (SHA-256) ---
Original File Hash: 323a1d7573ffa2eb128eb3196b892cfde06aa4c7fe52566a9ca6ce3f23c8c122
Expected Hash: 323a1d7573ffa2eb128eb3196b892cfde06aa4c7fe52566a9ca6ce3f23c8c122
Current Hash:  323a1d7573ffa2eb128eb3196b892cfde06aa4c7fe52566a9ca6ce3f23c8c122
[+] Integrity Verification Passed: File is untouched.

--- 4. ERROR HANDLING TESTS ---
Testing missing file:
[-] Error: [Errno 2] No such file or directory: 'non_existent_file.enc'

Testing corrupted file decryption:
[-] Error: Decryption failed. Invalid key or corrupted file.
\end{lstlisting}

\section{Network Traffic Filtering}
Traffic filtering rules were applied on the host environment using \texttt{iptables} to isolate guest networks and protect the server.

\subsection{Firewall Rules}
\begin{lstlisting}
iptables -A INPUT -m state --state ESTABLISHED,RELATED -j ACCEPT
iptables -A INPUT -p tcp -s 10.0.20.0/24 --dport 22 -j DROP
iptables -A INPUT -p tcp -s 10.0.10.0/24 --dport 22 -j ACCEPT
iptables -A INPUT -p tcp --dport 22 -j DROP
\end{lstlisting}

\subsection{Test Results}
\begin{itemize}
    \item \textbf{Staff Connection (10.0.10.15):} Permitted (\texttt{Connection Succeeded}).
    \item \textbf{Guest Connection (10.0.20.45):} Blocked (\texttt{Connection Timed Out}).
    \item \textbf{External Connection (172.16.0.88):} Blocked (\texttt{Connection Timed Out}).
\end{itemize}

\section{Conclusion and Submission}
The implementation secures student data both at rest (via AES encryption) and in transit (via SHA-256 hashing and host traffic filtering). Network rules effectively isolate unauthorized subnets from sensitive services.

\textbf{GitHub Repository URL:} \url{https://github.com/cynthianaissa7/cryptography-network-security-exam}

\end{document}
